% ==================== 5.1 引言 ====================
\section{引言}
\label{sec:ch5_intro}

月面岩石是巡视器自主行走面临的主要障碍之一。准确识别岩石目标、精确估计其位置和尺寸，并据此规划避障路径，是保障巡视器行走安全的关键。本章在数字孪生框架下，研究月面岩石识别与避障方法。

月面岩石识别面临以下挑战：

\textbf{（1）岩石外观多样}

月面岩石的形状、大小、颜色、纹理差异显著。从小于1厘米的碎块到数米的巨型岩石均有分布。岩石表面风化程度不同，呈现不同的视觉特征。

\textbf{（2）光照条件变化}

月面光照随太阳高度角变化而变化，同一岩石在不同光照条件下呈现不同的外观。暗影区域光照不足，岩石识别困难；高光照条件下，岩石与背景对比度下降。

\textbf{（3）月尘干扰}

月尘可能附着在岩石表面或相机镜头上，降低成像质量，影响识别精度。

\textbf{（4）实时性要求}

巡视器行走过程中需要实时识别前方岩石，为避障决策提供及时信息。

针对上述挑战，本章在数字孪生平台下构建岩石识别框架，开发增强型ORB-SLAM2系统实现岩石的精确识别与定位，并设计基于识别结果的避障策略。
